% from ega4-2-fr.tex
\oldpage[IV-2]{51}

paramètres pour $M$ $(0,16.3.6)$. Si le $A$-module
\[
N=M\Big/\Big(\sum_{i=1}^k x_iM\Big)
\]
est intègre (3.2.4), alors $M$ est intègre et la suite
$(x_i)_{1\leqslant i\leqslant k}$ est $M$-régulière.
\end{corollary}

\begin{proof}
Par récurrence sur $k$, on est aussitôt ramené au cas $k=1$~; nous écrirons
$x$ au lieu de $x_1$~; l'hypothèse que $x$ fait partie d'un système de
paramètres pour $M$ entraîne que
$\dim(N)=\dim(M)-1$ $(0,16.3.7)$. Posons $n=\dim(N)$~; il y a donc un élément
minimal $\mathfrak p$ de $\operatorname{Supp}(M)$ tel que
$\dim(M/\mathfrak pM)=n+1$ $(0,16.3.4)$, et pour tout entier $j>0$ on a alors
aussi $\dim(M/\mathfrak p^jM)=n+1$ $(0,16.3.5)$~; en outre $x$ fait partie
d'un système de paramètres pour $M/\mathfrak p^jM$ $(0,16.3.5)$, donc, si
l'on pose $M'=M/\mathfrak p^jM$ et $N'=M'/xM'$, on a $\dim(N')=n$. Il est
clair que l'on a un homomorphisme surjectif $v:N\to N'$~; montrons que $v$ est
bijectif. En effet, si $P=\operatorname{Ker}(v)$, on a
$\operatorname{Ass}(P)\subset\operatorname{Ass}(N)$ et puisque $N$ est
intègre, l'hypothèse $P\neq0$ entraînerait que $\operatorname{Ass}(P)$ et
$\operatorname{Ass}(N)$ seraient tous deux réduits à l'unique point
$\mathfrak q$ de
$\operatorname{Ass}(N)=\operatorname{Supp}(N)$~; mais comme
$\dim(N')=\dim(N)$, on a $N'_{\mathfrak q}\neq0$, et l'hypothèse
$\operatorname{long}(N_{\mathfrak q})=1$ entraîne
$\operatorname{long}(N'_{\mathfrak q})=1$ puisque
$N_{\mathfrak q}\to N'_{\mathfrak q}$ est surjective. On aurait donc
$P_{\mathfrak q}=0$, contrairement à l'hypothèse, d'où notre assertion.
Mais alors $N'$, étant isomorphe à $N$, est intègre~; en outre, le support de
$N'$ (égal à l'intersection de $\operatorname{Supp}(M')$ et de $V(x)$) ne
peut contenir $\operatorname{Supp}(M')$, et ce dernier ensemble est
irréductible par construction. L'hypothèse que $M'/xM'$ est intègre (donc
réduit) entraîne alors que $x$ est $M'$-régulier en vertu de (3.4.6). On en
conclut que le noyau de l'homothétie $z\mapsto xz$ dans $M$ est contenu dans
$\mathfrak p^jM\subset\mathfrak m^jM$ pour tout entier $j$, et ce noyau est
donc réduit à $0$ $(0_{\mathrm I},7.3.5)$, ce qui prouve que $x$ est
$M$-régulier. On peut alors appliquer (3.4.5), ce qui prouve que $M$ est
intègre.
\end{proof}

\begin{remark}[3.4.8]
\label{IV.2.3.4.8-fr}
La proposition analogue à (3.4.7), où on remplace «~intègre~» par
«~réduit~» n'est plus nécessairement exacte. Considérons par exemple l'anneau
de polynômes $C=K[X,Y,Z]$ sur un corps $K$, l'anneau quotient
$B=C/\mathfrak p\mathfrak q$, où $\mathfrak p=CZ$,
$\mathfrak q=CX^2+CY$~; soit $A$ l'anneau local de $B$ correspondant à l'idéal
maximal image de $CX+CY+CZ$ dans $B$. Si $x$, $z$ sont les images canoniques
de $X$, $Z$ dans $A$, il est clair que $xz\neq0$ mais $x^2z^2=0$~; d'autre
part, comme $A/xA$ est isomorphe à $K[Y,Z]/(YZ)$, on a
$\dim(A/xA)=1$ tandis que $\dim(A)=2$, donc $x$ appartient à un système de
paramètres de $A$ $(0,16.3.4)$, $A/xA$ est réduit, mais $A$ ne l'est pas.
\end{remark}

\begin{proposition}[3.4.9]
\label{IV.2.3.4.9-fr}
Soient $A$ un anneau noethérien, $M$ un $A$-module de type fini, $f$ un
élément $M$-régulier de $A$ tel que $M/fM$ n'ait pas d'idéaux premiers
associés immergés. Si $\mathfrak p_i$ $(1\leqslant i\leqslant m)$ sont les
idéaux premiers associés à $M/fM$, alors, pour tout entier $n>0$, $f^nM$ est
l'intersection des images réciproques des $f^nM_{\mathfrak p_i}$ par les
applications canoniques $M\to M_{\mathfrak p_i}$ $(1\leqslant i\leqslant m)$.
\end{proposition}

\begin{proof}
Tout revient à montrer que les $\mathfrak p_i$ sont aussi les idéaux premiers
associés à $M/f^nM$, car alors les saturés de $f^nM$ pour les
$\mathfrak p_i$ sont les sous-modules de la décomposition primaire réduite
(nécessairement unique) de $f^nM$ dans $M$. Or, on a
\[
\operatorname{Ass}(f^{n-1}M/f^nM)
\subset\operatorname{Ass}(M/f^nM)
\subset\operatorname{Ass}(M/f^{n-1}M)
\cup\operatorname{Ass}(f^{n-1}M/f^nM)
\]
par (3.1.7), et comme $f$ est $M$-régulier,
$f^{n-1}M/f^nM$ est isomorphe à $M/fM$~; il suffit alors de raisonner par
récurrence sur $n$.
\end{proof}
